\documentclass[aip,jcp,reprint,noshowkeys,superscriptaddress]{revtex4-1}
\usepackage{graphicx,dcolumn,bm,xcolor,microtype,multirow,amscd,amsmath,amssymb,amsfonts,physics,longtable,wrapfig}
\usepackage[version=4]{mhchem}

\usepackage[utf8]{inputenc}
\usepackage[T1]{fontenc}
\usepackage{txfonts}

\usepackage[
    colorlinks,
    linkcolor={red!50!black},
    citecolor={red!70!black},
    urlcolor={red!80!black}
	]{hyperref}
\urlstyle{same}

\newcommand{\ie}{\textit{i.e.}}
\newcommand{\eg}{\textit{e.g.}}
\newcommand{\alert}[1]{\textcolor{red}{#1}}
\usepackage[normalem]{ulem}
\newcommand{\titou}[1]{\textcolor{red}{#1}}
\newcommand{\trashPFL}[1]{\textcolor{\red}{\sout{#1}}}
\newcommand{\PFL}[1]{\titou{(\underline{\bf PFL}: #1)}}

\newcommand{\cre}[1]{\Hat{a}_{#1}^{\dag}}
\newcommand{\ani}[1]{\Hat{a}_{#1}}
\newcommand{\ca}[2]{\Hat{a}^{#1}_{#2}}

\newcommand{\mc}{\multicolumn}
\newcommand{\fnm}{\footnotemark}
\newcommand{\fnt}{\footnotetext}
\newcommand{\tabc}[1]{\multicolumn{1}{c}{#1}}
\newcommand{\QP}{\textsc{quantum package}}
\newcommand{\T}[1]{#1^{\intercal}}

% coordinates
\newcommand{\br}{\boldsymbol{r}}
\newcommand{\bx}{\boldsymbol{x}}
\newcommand{\bI}{\boldsymbol{1}}
\newcommand{\cK}{\mathcal{K}}
\newcommand{\bG}{\boldsymbol{G}}
\newcommand{\bH}{\boldsymbol{H}}
\newcommand{\bM}{\boldsymbol{M}}
\newcommand{\bE}{\boldsymbol{E}}
\newcommand{\bSigma}{\boldsymbol{\Sigma}}

% orbital energies
\newcommand{\eps}{\epsilon}
\newcommand{\irr}{\text{irr}}

% shortcuts for greek letters
\newcommand{\si}{\sigma}
\newcommand{\la}{\lambda}
\newcommand{\ii}{\mathrm{i}}
\newcommand{\up}{\uparrow}
\newcommand{\dw}{\downarrow}
\newcommand{\upup}{\up\up}
\newcommand{\updw}{\up\dw}
\newcommand{\dwup}{\dw\up}
\newcommand{\dwdw}{\dw\downarrow}

\newcommand{\h}{\text{1h}}
\newcommand{\p}{\text{1p}}
\newcommand{\hp}{\text{1h+1p}}
\newcommand{\hhp}{\text{2h1p}}
\newcommand{\pph}{\text{2p1h}}
\newcommand{\ph}{\text{ph}}
\newcommand{\eh}{\text{eh}}
\newcommand{\he}{\text{he}}
\newcommand{\ee}{\text{ee}}
\newcommand{\teh}{\overline{\text{eh}}}
\newcommand{\pp}{\text{pp}}
\newcommand{\hh}{\text{hh}}
\newcommand{\xc}{\text{xc}}


% addresses
\newcommand{\LCPQ}{Laboratoire de Chimie et Physique Quantiques (UMR 5626), Universit\'e de Toulouse, CNRS, UPS, France}


% statement-status markers used throughout the notes
\newcommand{\Exact}{\textbf{Exact identity.}\ }
\newcommand{\Established}{\textbf{Established result.}\ }
\newcommand{\Proposed}{\textbf{Proposed approximation.}\ }

\begin{document}

\title{Frequency dependence in EOM theories and its relation to the BSE}

%\author{Pierre-Fran\c{c}ois \surname{Loos}}
%	\email{loos@irsamc.ups-tlse.fr}
%	\affiliation{\LCPQ}
\iffalse
\begin{abstract}
bla bla bla
\end{abstract}
\fi

\maketitle


\section{Equation-of-Motion Theories}

We start by analyzing the frequency dependence in Equation-of-Motion (EOM) theories.
Neutral excited-states within a EOM approach are determined from 
\begin{align}
 \langle \Psi_0 | \hat{C}^\dagger_\mu [\hat{H}, \hat{R}_n] | \Psi_0 \rangle = \omega_n \langle \Psi_0 | \hat{C}^\dagger_\mu \hat{R}_n | \Psi_0 \rangle,
 \label{eq:EOM}
\end{align}
with $\hat{R}_n$ denoting the excitation operator for the $n$-th excited state 
\begin{align}
  \hat{R}_n = \sum_\mu r_\mu^{(n)} \hat{C}_\mu.
\end{align}
Eq.~\eqref{eq:EOM} can be written in matrix form as
\begin{align}
 \mathbf{H} \mathbf{r}_n = \omega_n  \mathbf{r}_n.
\end{align}
Assuming the excitation operator to be restricted to single and double excitations, i.e.,
\begin{align}
  \hat{R}_n = \sum_{ia} r_i^{a,(n)} \hat{a}^\dagger_a \hat{a}_i + \frac{1}{4} \sum_{ijab} r_{ij}^{ab,(n)} \hat{a}^\dagger_a \hat{a}^\dagger_b \hat{a}_j \hat{a}_i,
\end{align}
where $i,j$ ($a,b$) label occupied (virtual) orbitals.
The block-structure of the Hamiltonian matrix $\mathbf{H}$ in the basis of single and double excitations is then given by
\begin{align}
\mathbf{H} = \begin{pmatrix}
 \mathbf{H}_{\text{ss}} & \mathbf{H}_{\text{sd}} \\
 \mathbf{H}_{\text{ds}} & \mathbf{H}_{\text{dd}}
\end{pmatrix},
\end{align}
where the subscripts `s' and `d' refer to single and double excitations, respectively.
Downfolding the Hamiltonian in the space of single excitations leads to the frequency-dependent effective Hamiltonian
\begin{align}
\mathbf{H}_{\text{eff}}(\omega) = \mathbf{H}_{\text{ss}} + \mathbf{H}_{\text{sd}} (\omega - \mathbf{H}_{\text{dd}})^{-1} \mathbf{H}_{\text{ds}}.
\end{align}

For the following analysis we assume a diagonal doubles block, i.e.,
\begin{align}
\mathbf{H}_{\text{dd}} = \langle \Phi_{ij}^{ab} | \hat{H} | \Phi_{kl}^{cd} \rangle = \delta_{ik} \delta_{jl} \delta_{ac} \delta_{bd} (\epsilon_a + \epsilon_b - \epsilon_i - \epsilon_j)
\end{align}
where $\epsilon_p$ denote the canonical Hartree--Fock orbital energies.
The coupling block read
\begin{align}
  \mathbf{H}_{\text{sd}} &= \langle \Phi_i^a | \hat{H} | \Phi_{jk}^{bc} \rangle  \nonumber \\
  &= \delta_{ik} \langle ab||jc \rangle - \delta_{ij} \langle ab||kc \rangle\nonumber \\
  &+ \delta_{ab} \langle jk||ic \rangle - \delta_{ac} \langle jk||ib \rangle .
\end{align}
This illustrates that we can distinguish between to fundamental building blocks in the frequency-dependent part of the effective Hamiltonian.

These can be categorized into two types:
Genuine two-particle interaction terms
\begin{align}
   \Gamma_{iajb}(\omega) = \sum_{kc} \frac{\langle ab||kc \rangle \langle kc||ij \rangle}{\omega - (\epsilon_c + \epsilon_b - \epsilon_k - \epsilon_i)},
\end{align}
or terms with a spectator particle 
\begin{align}
  \delta_{ab} \Sigma_{jl}(\omega') \nonumber \\
  &= \delta_{ab} \sum_{klc} \frac{\langle jk||ic \rangle \langle ic||kl \rangle}{\omega - (\epsilon_c + \epsilon_b - \epsilon_i - \epsilon_k)} \nonumber \\
  &= \delta_{ab} \Sigma_{jl}(\omega - \epsilon_b).
\end{align}
The latter can be directly identified as self-energy-like term dressing the single-particle propagator.

With this, we can rewrite the frequency-dependent effective Hamiltonian as a sum of genuine two-particle interaction terms and disconnected one-particle contributions
\begin{align}
\mathbf{H}_{\text{eff}}(\omega)_{ijab} &= \delta_{ij} \left( \delta_{ab} \epsilon_a + \Sigma_{ab}(\omega + \epsilon_j) \right) \nonumber \\
&- \delta_{ab} \left(\delta_{ij} \epsilon_i +  \Sigma_{ij}(\omega - \epsilon_b) \right) + V_{iajb} + \Gamma_{iajb}(\omega).
\end{align}
\textcolor{red}{This is just for illustration purposes. The signs might be wrong.}

Summarizing the frequency dependence as $\nu = \omega - \epsilon_b$ and $\nu' = \omega - \epsilon_j$, we can rewrite the effective Hamiltonian in terms of these shifted frequencies
\begin{align}
\mathbf{H}_{\text{eff}}(\omega, \nu, \nu')_{ijab} &= \delta_{ij} \left( \delta_{ab} \epsilon_a + \Sigma_{ab}(\nu') \right) \nonumber \\
&- \delta_{ab} \left(\delta_{ij} \epsilon_i +  \Sigma_{ij}(\nu) \right) + V_{iajb} + \Gamma_{iajb}(\omega).
\end{align}


\textbf{Interpretation:}
The diagrammatic rules for constructing the BSE, i.e., via the Parquet approximation, do not allow for the inclusion of the self-energy insertions directly, due to two-particle reducible and irreducible contributions being separated from one-particle contributions to avoid double counting from the usage of a interacting one-particle Green's function.
The fermionic frequencies therefore reflect the fact that the evaluation of the self-energy has to be shifted by the corresponding bosonic frequency compared to the single-particle treatment via the Dyson equation. 

\textbf{Next:}
Comparison with the multifrequency kernels derived by Titou.
Furthermore, the derivation above decouples excitations from deexcitations. 
I conjecture that this is the reason why the particle/hole self-energy insertions appear only for 2h1p/2p1h intermediate states.
In the CC case, one would have additional static 2p1h/2h1p contributions arising from the contraction with the $T_2$ amplitudes.

%%%%%%%%%%%%%%%%%%%%%%%%
\bibliography{parquet-mr-hedin}
%%%%%%%%%%%%%%%%%%%%%%%%

\end{document}

